%!TEX root = ../dokumentation.tex

\chapter{Messbarer Organisationaler Nutzen (\ac{MOV})}
\label{chap:mov}

\section{\ac{MOV}-Statement}
\label{sec:mov-statement}

Das Projekt „Lernplattform" ist erfolgreich, wenn es innerhalb von 12 Monaten
nach Projektabschluss den administrativen Aufwand im Schulungsmanagement um
mindestens 30\,\% (~3\,Mio.\,€\,p.\,a.) reduziert, eine aktive Systemnutzung
von 60\,\% der Mitarbeitenden erreicht und aggregierte Kompetenzkennzahlen als
Basis für organisationsweite Entwicklungs- und Bonimodelle bereitstellt.

\section{Wirkungsbereiche und Wertbeitrag}
\label{sec:mov-wirkungsbereiche}

Folgende Wirkungsbereiche werden adressiert:

\begin{itemize}
    \item \textbf{Finanziell (Cheaper):} Reduktion der Prozesskosten durch
    Ablösung manueller, heterogener Verwaltungstätigkeiten.

    \item \textbf{Operativ (Better/Faster):} Harmonisierung der End-to-End-Prozesse
    (Planung, Buchung, Dokumentation) zur Steigerung der administrativen Effizienz.

    \item \textbf{Strategisch (Do More):} Steigerung der Innovationsfähigkeit durch
    kontinuierliche Kompetenzentwicklung und Stärkung der Arbeitgebermarke
    (Employer~Branding).

    \item \textbf{Mitarbeitende (Social):} Etablierung einer Kultur des lebenslangen
    Lernens durch transparenten Zugang zu Entwicklungsperspektiven.
\end{itemize}

\section{Metriken und Zielwerte}
\label{sec:mov-metriken}

Die Überprüfung des Projekterfolgs erfolgt anhand der definierten Metriken.

\begin{xltabular}{\textwidth}{@{}lXXX@{}}
    \caption{\ac{MOV}-Metriken und Zielwerte} \label{tab:mov-metriken} \\
    \hline
    \textbf{Dimension} & \textbf{Metrik / Kennzahl} & \textbf{Zielwert} & \textbf{Zeitraum (nach Go-Live)} \\
    \hline
    \endfirsthead

    \hline
    \textbf{Dimension} & \textbf{Metrik / Kennzahl} & \textbf{Zielwert} & \textbf{Zeitraum (nach Go-Live)} \\
    \hline
    \endhead

    \hline
    \endfoot

    Finanziell &
    Reduktion administrativer Aufwand &
    $\geq$\,30\,\% (ca.\,3\,Mio.\,€) &
    12 Monate \\

    Operativ &
    Systemnutzungsquote (Aktiv) &
    $\geq$\,60\,\% der Mitarbeitenden &
    12 Monate \\

    Strategisch &
    \ac{KPI}-Verfügbarkeit &
    Power\,BI-Dashboards aktiv &
    Bei Projektabschluss \\

    Qualitativ &
    Datenintegrität &
    \textit{Single Source of Truth} aktiv &
    Bei Projektabschluss \\

\end{xltabular}

\section{Zeitrahmen und Verifikation}
\label{sec:mov-verifikation}

Der volle organisationale Nutzen wird innerhalb von 12 Monaten nach
Systemübergabe realisiert. Die Verifikation erfolgt durch administrative
Analyse (Personalbindungskosten im Schulungsmanagement vor und nach
Einführung), System-Audits (Log-Daten für Nutzungsquote, \ac{MOV}\,2) sowie
Reporting-Check (Dashboard-Funktionalität für Kompetenzkennzahlen durch
\ac{HR}-Leitung).
